\chapter{Introduzione}\label{ch:introduzione}

\section{Descrizione}
\textbf{DriveHub} è un'applicazione desktop che riunisce le attività essenziali di una concessionaria: consultazione del catalogo, vendita e noleggio di veicoli, prenotazione di test drive e valutazione dei veicoli proposti dai clienti. L'obiettivo è collegare queste attività alle stesse informazioni su veicoli, disponibilità e operazioni economiche, evitando di trattarle come registrazioni indipendenti.

La terminologia segue il punto di vista della concessionaria. Una \emph{vendita al cliente} trasferisce un veicolo dalla concessionaria al cliente; un \emph{acquisto dal cliente}, detto anche \emph{cessione del cliente}, riguarda invece un veicolo proposto dal cliente alla concessionaria. Le richieste in attesa di esito sono distinte dalle operazioni registrate nel registro economico.

Il sistema prevede tre ruoli, con aree applicative dedicate:
\begin{itemize}
  \item \textbf{Cliente:} si registra, consulta i cataloghi dei veicoli in vendita e a noleggio, prenota test drive e noleggi e propone un proprio veicolo alla concessionaria. Per avviare una vendita al cliente versa un acconto; dopo l'approvazione del venditore può pagare il saldo. Consulta le proprie richieste e prenotazioni e gestisce il conto dimostrativo e il profilo.
  \item \textbf{Venditore:} esamina le richieste di vendita al cliente e le proposte di cessione, approvandole o rifiutandole. Gestisce inoltre la destinazione dei veicoli tra vendita e noleggio.
  \item \textbf{Manager:} gestisce marche, modelli e veicoli, sceglie la base del prezzo e applica sconti a modelli o singoli veicoli. Consulta il registro economico, che distingue vendite al cliente, acquisti dal cliente e penali.
\end{itemize}

La disponibilità dipende dallo stato del veicolo e dalle operazioni che lo impegnano. Per esempio, il versamento dell'acconto blocca un veicolo in vendita; il completamento della vendita lo porta allo stato venduto. Noleggi e test drive richiedono anche controlli sui periodi di prenotazione. Le regole sono coordinate dai servizi applicativi e sostenute dai vincoli del database.

\clearpage
\section{Stack tecnologico}
Il progetto usa le seguenti tecnologie, riscontrate nelle dipendenze e nei sorgenti correnti:
\begin{itemize}
  \item \textbf{Java 21:} implementa dominio, servizi applicativi, accesso ai dati e presentazione. Il progetto Maven imposta a 21 sia la versione dei sorgenti sia quella di destinazione del compilatore.
  \item \textbf{JavaFX 21.0.2:} realizza l'interfaccia desktop con controlli e layout costruiti direttamente in Java e uno stile CSS condiviso. Le viste invocano i servizi applicativi; le tabelle usano caricamenti asincroni per le letture previste, con aggiornamento dei controlli sul thread JavaFX.
  \item \textbf{JDBC e PostgreSQL 16:} i DAO eseguono SQL e mappano i risultati negli oggetti di dominio, senza un ORM. Il driver PostgreSQL dichiarato è la versione 42.7.3. Il gestore delle connessioni offre transazioni con commit e rollback; lo schema usa vincoli relazionali e tipi enumerati nativi per stati, tipi di richiesta, modalità del prezzo e tipi di registrazione economica.
  \item \textbf{Maven:} gestisce dipendenze, compilazione, test e avvio tramite il plugin JavaFX. Gli script Windows forniscono un percorso di avvio con PostgreSQL 16 in Docker; è previsto anche l'uso di un server PostgreSQL installato separatamente.
  \item \textbf{JUnit 5:} le dipendenze di test dichiarano la versione 5.10.2. Il progetto distingue i test senza database dalle prove di integrazione PostgreSQL, attivate con un profilo dedicato e credenziali di test separate.
\end{itemize}

La presentazione, i servizi, gli oggetti di dominio e i DAO hanno responsabilità distinte. \code{AppContext} costruisce e collega esplicitamente le dipendenze, mentre \code{App} avvia JavaFX e apre la dashboard del ruolo autenticato. Questa organizzazione non richiede un framework di dependency injection o un servizio web intermedio: l'applicazione desktop accede al database tramite JDBC.

\clearpage
\section{Semplificazioni adottate}\label{sec:semplificazioni}
DriveHub è un progetto didattico: rappresenta i flussi necessari a dimostrare le regole applicative, con i seguenti limiti di perimetro.
\begin{itemize}
  \item \textbf{Conto dimostrativo.} Il conto associato al cliente conserva il saldo nel database. Ricariche, addebiti, accrediti e rimborsi sono movimenti interni all'applicazione: non intervengono banche o circuiti di pagamento e l'IBAN generato non prova l'esistenza di un conto reale.
  \item \textbf{Registrazione e personale.} La schermata pubblica e il servizio di registrazione creano soltanto clienti. Venditori e manager devono essere già presenti nel database; non è implementato un percorso applicativo per amministrare il personale. I dati dimostrativi comprendono utenti per i tre ruoli.
  \item \textbf{Regole economiche circoscritte.} Prezzi e tariffe sono calcolati dalle regole del progetto. La base manuale scelta dal manager sostituisce quella del modello nella stessa formula, conservando gli effetti di sconti e svalutazioni: non impone un prezzo finale fisso. Il totale della vendita al cliente viene concordato al versamento dell'acconto del 20\%; il saldo usa quel totale, anche se in seguito cambiano basi o sconti.
  \item \textbf{Ambiente di esecuzione.} La persistenza richiede un database PostgreSQL compatibile, già predisposto. L'applicazione Java non crea lo schema né carica dati di esempio all'avvio e non dispone di un database alternativo in memoria. L'inizializzazione dello schema e dei dati dimostrativi è un'attività separata; lo script Windows la esegue soltanto nella prima configurazione del proprio database Docker. I dati persistenti non vengono ricreati a ogni avvio.
  \item \textbf{Processi aziendali ridotti.} Il modello copre catalogo, richieste, prenotazioni e registrazioni economiche. Non modella una rete di sedi, fornitori, officina, consegna fisica, gestione documentale dei contratti, fatturazione o assistenza successiva alla vendita.
\end{itemize}

Queste scelte delimitano ciò che il prototipo dimostra. Il conto e i dati di esempio rendono riproducibili i flussi, ma non equivalgono a un'integrazione bancaria o a un sistema completo per l'esercizio di una concessionaria.
